\fronttocchapter{索引}
\label{back:index}
\chaptercredit{}

索引按英文词项排列，词项后的数字是原书印刷页码。点击页码可跳到该页附近的译文章节，再用中英文词项搜索具体位置。下表列出各章和附录在原书与译本中的起始页。

\begin{center}
\small\renewcommand{\arraystretch}{1.08}
\begin{tabularx}{\linewidth}{@{}Xrr@{}}
\toprule
章节 & 原版起页 & 译本起页\\
\midrule
1\quad 大语言模型导论 & 1 & \pageref{ch:intro}\\
2\quad 语言模型预训练 & 29 & \pageref{ch:pretraining}\\
3\quad 基于提示的学习 & 83 & \pageref{ch:prompting}\\
4\quad 大语言模型的适配与使用 & 135 & \pageref{ch:adaptation}\\
5\quad 大语言模型的对齐调优 & 177 & \pageref{ch:alignment}\\
6\quad 大语言模型的挑战与应对方法 & 219 & \pageref{ch:challenges}\\
7\quad 检索增强生成 & 275 & \pageref{ch:rag}\\
8\quad 生产环境中的大语言模型 & 315 & \pageref{ch:production}\\
9\quad 多模态大语言模型 & 375 & \pageref{ch:multimodal}\\
10\quad 大语言模型的演进与前沿 & 423 & \pageref{ch:frontiers}\\
附录 A\quad 深度学习基础 & 439 & \pageref{app:deep-learning}\\
附录 B\quad 强化学习基础 & 451 & \pageref{app:reinforcement-learning}\\
\bottomrule
\end{tabularx}
\end{center}

\begingroup
\newcommand{\bookindexgroup}[1]{\par\medskip\noindent\textsf{\textbf{#1}}\par\nobreak\smallskip}
\newcommand{\bookindexpage}[2]{\hyperlink{#1}{\textcolor{blue!50!black}{#2}}}
\newcommand{\bookindexentry}[3]{%
  \par\hangindent=1.25em\hangafter=1\noindent
  \if\relax\detokenize{#1}\relax
    \textbf{#2}%
  \else
    \textbf{#1}\hskip.25em\relax(#2)%
  \fi
  \enspace#3\par}
\small
\raggedright
\setlength{\parskip}{2pt}
\setlength{\columnsep}{12pt}
\interlinepenalty=10000
\begin{multicols}{2}
% 对应原书印刷页 465--472。数字是原书定位页码，链接目标是译文章节。
\bookindexgroup{0–9}
\bookindexentry{三维物体}{3D Objects}{\bookindexpage{section.9.3}{378}}
\bookindexgroup{A}
\bookindexentry{}{AdaLoRA}{\bookindexpage{section.4.3}{152}}
\bookindexentry{Adam 优化器}{Adam Optimizer}{\bookindexpage{section.2.5}{58}}
\bookindexentry{AdamW 优化器}{AdamW Optimizer}{\bookindexpage{section.2.5}{58}}
\bookindexentry{适应性预训练}{Adaptive Pre-training}{\bookindexpage{section.1.5}{15–17}, \bookindexpage{section.8.6}{343}}
\bookindexentry{对抗攻击}{Adversarial Attacks}{\bookindexpage{section.6.3}{248}}
\bookindexentry{智能体}{Agents}{\bookindexpage{section.5.3}{188}}
\bookindexentry{人工智能助手}{AI Assistants}{\bookindexpage{section.8.2}{317}}
\bookindexentry{}{ALBERT}{\bookindexpage{section.2.6}{64}, \bookindexpage{section.3.3}{112}}
\bookindexentry{对齐}{Alignment}{\bookindexpage{section.6.1}{222}}
\bookindexentry{对齐调优}{Alignment Tuning}{\bookindexpage{section.1.5}{18–19}, \bookindexpage{section.5.1}{178}–\bookindexpage{section.5.2}{184}, \bookindexpage{section.8.8}{358}}
\bookindexentry{}{AlignScore}{\bookindexpage{section.6.1}{222}}
\bookindexentry{}{Alpaca}{\bookindexpage{section.10.2}{424}}
\bookindexentry{交替语言建模}{Alternate Language Modeling}{\bookindexpage{section.2.5}{53}}
\bookindexentry{答案工程}{Answer Engineering}{\bookindexpage{section.3.4}{113–115}}
\bookindexentry{答案忠实度}{Answer Faithfulness}{\bookindexpage{section.7.5}{294}, \bookindexpage{section.7.5}{297–300}}
\bookindexentry{答案映射}{Answer Mapping}{\bookindexpage{section.3.2}{94}}
\bookindexentry{答案相关性}{Answer Relevance}{\bookindexpage{section.7.5}{294}, \bookindexpage{section.7.5}{300–301}}
\bookindexentry{答案搜索}{Answer Search}{\bookindexpage{section.3.2}{93–94}}
\bookindexentry{}{Anthropic}{\bookindexpage{section.1.4}{11}, \bookindexpage{section.5.5}{206}, \bookindexpage{section.5.5}{212}, \bookindexpage{section.5.6}{214}, \bookindexpage{section.8.5}{341}, \bookindexpage{section.10.2}{424}}
\bookindexentry{}{ArXiv}{\bookindexpage{section.9.5}{400}}
\bookindexentry{注意力}{Attention}{\bookindexpage{section.1.3}{9}, \bookindexpage{section.1.4}{11}, \bookindexpage{section.1.5}{15}, \bookindexpage{section.2.2}{32}–\bookindexpage{section.2.3}{33}, \bookindexpage{section.2.3}{41}, \bookindexpage{section.8.5}{339}}
\bookindexentry{属性条件控制}{Attribute Conditioning}{\bookindexpage{section.6.3}{251}}
\bookindexentry{音频—文本模态}{Audio-Text Modality}{\bookindexpage{section.9.3}{384}}
\bookindexentry{}{AudioLDM}{\bookindexpage{section.9.6}{411}}
\bookindexentry{}{AutoGPT}{\bookindexpage{section.10.4}{433}}
\bookindexentry{自动答案映射}{Automated Answer Mapping}{\bookindexpage{section.3.4}{115}}
\bookindexentry{自动构造提示}{Automated Prompt}{\bookindexpage{section.3.3}{103}}
\bookindexentry{自动驾驶导航}{Autonomous Vehicle Navigation}{\bookindexpage{section.9.3}{379}}
\bookindexentry{}{AutoPrompt}{\bookindexpage{section.3.3}{106}, \bookindexpage{section.3.3}{110}}
\bookindexentry{自回归解码}{Autoregressive Decoding}{\bookindexpage{section.2.5}{58}}
\bookindexentry{}{AWS Bedrock}{\bookindexpage{section.8.6}{346}, \bookindexpage{section.8.7}{351}}
\bookindexentry{}{AWS Sagemaker Jumpstart}{\bookindexpage{section.8.6}{346}}
\bookindexgroup{B}
\bookindexentry{}{BabelNet}{\bookindexpage{section.6.2}{231}}
\bookindexentry{}{BabyAGI}{\bookindexpage{section.10.4}{433}}
\bookindexentry{反向传播}{Backpropagation}{\bookindexpage{section.5.2}{185}, \bookindexpage{section.6.3}{253}}
\bookindexentry{}{BART}{\bookindexpage{section.3.3}{110}}
\bookindexentry{批归一化}{BatchNorm}{\bookindexpage{section.2.3}{39}}
\bookindexentry{贝叶斯网络}{Bayesian Networks}{\bookindexpage{section.1.3}{8}}
\bookindexentry{束搜索}{Beam Search}{\bookindexpage{section.2.5}{59}, \bookindexpage{section.5.5}{210}}
\bookindexentry{}{BEATs}{\bookindexpage{section.9.3}{379}}
\bookindexentry{}{BentoML}{\bookindexpage{section.10.3}{430}}
\bookindexentry{}{BERT}{\bookindexpage{section.1.4}{10}, \bookindexpage{section.2.4}{43}, \bookindexpage{section.2.6}{61–64}, \bookindexpage{section.2.7}{76}, \bookindexpage{section.3.1}{86}, \bookindexpage{section.8.3}{325}}
\bookindexentry{}{BERTScore}{\bookindexpage{section.8.3}{323}, \bookindexpage{section.8.3}{325–326}}
\bookindexentry{}{BGE}{\bookindexpage{section.7.4}{283}}
\bookindexentry{偏见}{Bias}{\bookindexpage{section.6.2}{229}–\bookindexpage{section.6.3}{247}}
\bookindexentry{偏见缓解}{Bias Mitigation}{\bookindexpage{section.6.2}{241}}
\bookindexentry{BIG-bench 研讨会}{Big-bench workshop}{\bookindexpage{section.5.4}{190}}
\bookindexentry{}{BigScience}{\bookindexpage{section.8.2}{321}}
\bookindexentry{}{Bingo}{\bookindexpage{section.9.4}{399}}
\bookindexentry{}{BioMistral}{\bookindexpage{section.8.3}{322}}
\bookindexentry{}{BitFit}{\bookindexpage{section.4.3}{154}}
\bookindexentry{位宽}{Bitwidth}{\bookindexpage{section.4.4}{158}}
\bookindexentry{黑箱探查}{Black-box Probing}{\bookindexpage{section.6.4}{258}}
\bookindexentry{}{BLEU}{\bookindexpage{section.8.3}{323–324}}
\bookindexentry{屏蔽词表}{Blocklisting}{\bookindexpage{section.6.3}{252}}
\bookindexentry{}{BLOOM}{\bookindexpage{section.4.4}{159}, \bookindexpage{section.8.2}{321}, \bookindexpage{section.10.3}{430}}
\bookindexentry{}{BloombergGPT}{\bookindexpage{section.1.5}{17}, \bookindexpage{section.10.4}{431}}
\bookindexentry{}{BOLD}{\bookindexpage{section.6.2}{236}, \bookindexpage{section.6.2}{238}}
\bookindexentry{Bradley–Terry 模型}{Bradley-Terry Model}{\bookindexpage{section.5.4}{195}, \bookindexpage{section.5.5}{211}}
\bookindexentry{广度优先搜索}{Breadth-first Search}{\bookindexpage{section.4.6}{168}}
\bookindexentry{长度惩罚}{Brevity Penalty}{\bookindexpage{section.8.3}{324}}
\bookindexentry{}{BriVL}{\bookindexpage{section.9.3}{381}}
\bookindexgroup{C}
\bookindexentry{}{C-Former}{\bookindexpage{section.9.3}{379}}
\bookindexentry{}{Cassandra}{\bookindexpage{section.8.6}{346}}
\bookindexentry{灾难性遗忘}{Catastrophic Forgetting}{\bookindexpage{section.6.2}{244}}
\bookindexentry{类别偏见分数}{Categorical Bias Score}{\bookindexpage{section.6.2}{235}}
\bookindexentry{因果解码器架构}{Causal Decoder Architecture}{\bookindexpage{section.1.5}{16}, \bookindexpage{section.2.5}{55}}
\bookindexentry{后见链}{Chain-of-Hindsight}{\bookindexpage{section.5.1}{183}}
\bookindexentry{思维链}{Chain-of-Thought}{\bookindexpage{section.1.5}{21}, \bookindexpage{section.4.5}{166–167}, \bookindexpage{section.8.3}{328}}
\bookindexentry{聊天机器人}{Chatbots}{\bookindexpage{section.1.5}{22}, \bookindexpage{section.2.4}{47}, \bookindexpage{section.5.1}{180}, \bookindexpage{section.5.4}{189}, \bookindexpage{section.5.6}{213}, \bookindexpage{section.7.4}{281}, \bookindexpage{section.8.2}{316–317}, \bookindexpage{section.9.2}{376}}
\bookindexentry{}{ChatGPT}{\bookindexpage{section.1.1}{1–2}, \bookindexpage{section.1.4}{10–12}, \bookindexpage{section.4.2}{137}, \bookindexpage{section.4.4}{161}, \bookindexpage{section.6.2}{232}, \bookindexpage{section.6.3}{248}, \bookindexpage{section.6.4}{257}, \bookindexpage{section.7.2}{276}, \bookindexpage{section.8.2}{318}, \bookindexpage{section.8.5}{342}, \bookindexpage{section.8.8}{358}}
\bookindexentry{国际象棋}{Chess}{\bookindexpage{section.5.5}{209}}
\bookindexentry{}{Chinchilla}{\bookindexpage{section.2.3}{39}, \bookindexpage{section.2.4}{47}, \bookindexpage{section.2.6}{72}, \bookindexpage{section.9.5}{404}}
\bookindexentry{}{Chroma}{\bookindexpage{section.8.6}{346}}
\bookindexentry{}{CIFAR-10}{\bookindexpage{section.10.3}{428}}
\bookindexentry{}{Claude}{\bookindexpage{section.7.3}{280}}
\bookindexentry{}{Claude 2}{\bookindexpage{section.4.2}{137}}
\bookindexentry{}{Claude 2.1}{\bookindexpage{section.10.2}{424}}
\bookindexentry{}{Claude 3}{\bookindexpage{section.2.6}{72}}
\bookindexentry{}{ClinicalBERT}{\bookindexpage{section.2.5}{49}}
\bookindexentry{}{CLIP}{\bookindexpage{section.9.3}{377}}
\bookindexentry{}{CLIP ViT}{\bookindexpage{section.9.3}{378}}
\bookindexentry{点云模态编码}{Cloud Modality Encoding}{\bookindexpage{section.9.3}{379}}
\bookindexentry{云服务}{Cloud Services}{\bookindexpage{section.8.6}{346}}
\bookindexentry{完形填空式提示}{Cloze Prompt}{\bookindexpage{section.3.3}{101}–\bookindexpage{section.3.4}{114}}
\bookindexentry{}{CMMU}{\bookindexpage{section.9.4}{398}}
\bookindexentry{卷积神经网络}{CNNs, Convolutional Neural Networks}{\bookindexpage{section.1.3}{8}, \bookindexpage{section.6.3}{249}}
\bookindexentry{共现偏见分数}{Co-Occurrence Bias Score}{\bookindexpage{section.6.2}{236}}
\bookindexentry{代码自动补全}{Code Auto-completion}{\bookindexpage{section.8.2}{319}}
\bookindexentry{语码转换}{Code-Switching}{\bookindexpage{section.8.2}{321}}
\bookindexentry{}{Codex-Med}{\bookindexpage{section.10.4}{431}}
\bookindexentry{编程}{Coding}{\bookindexpage{section.8.2}{318}}
\bookindexentry{}{Command R}{\bookindexpage{section.2.6}{73}}
\bookindexentry{虚构}{Confabulation}{\bookindexpage{section.4.5}{167}}
\bookindexentry{}{ConfAIde}{\bookindexpage{section.6.4}{259}}
\bookindexentry{宪法式人工智能}{Constitutional AI}{\bookindexpage{section.5.5}{206}, \bookindexpage{section.5.5}{207}}
\bookindexentry{内容创作}{Content Creation}{\bookindexpage{section.8.2}{317–318}}
\bookindexentry{上下文关联检验}{Context Association Test}{\bookindexpage{section.6.2}{236}}
\bookindexentry{上下文整合}{Context Consolidation}{\bookindexpage{section.7.4}{290–291}}
\bookindexentry{上下文长度}{Context Length}{\bookindexpage{section.8.1}{315}}
\bookindexentry{上下文相关性}{Context Relevance}{\bookindexpage{section.7.5}{294–297}}
\bookindexentry{上下文窗口}{Context Window}{\bookindexpage{section.10.2}{425}}
\bookindexentry{上下文嵌入}{Contextual Embeddings}{\bookindexpage{section.6.2}{234}, \bookindexpage{section.8.3}{325}}
\bookindexentry{上下文嵌入关联检验}{Contextualized Embedding Association Test}{\bookindexpage{section.6.2}{234}}
\bookindexentry{连续批处理}{Continuous Batching}{\bookindexpage{section.8.7}{352}}
\bookindexentry{连续积分发放对齐机制}{Continuous Integrate-and-Fire Alignment Mechanism}{\bookindexpage{section.9.3}{379}}
\bookindexentry{连续提示}{Continuous Prompt}{\bookindexpage{section.3.3}{108–111}, \bookindexpage{section.6.2}{242}}
\bookindexentry{对比学习}{Contrastive Learning}{\bookindexpage{section.9.3}{385–386}}
\bookindexentry{对话式人工智能}{Conversational AI}{\bookindexpage{section.8.2}{317}}
\bookindexentry{余弦相似度}{Cosine Similarity}{\bookindexpage{section.7.4}{285}, \bookindexpage{section.7.5}{301}}
\bookindexentry{反事实数据增强}{Counterfactual Data Augmentation}{\bookindexpage{section.6.2}{238}, \bookindexpage{section.6.2}{241}}
\bookindexentry{反事实数据替换}{Counterfactual Data Substitution}{\bookindexpage{section.6.2}{241}}
\bookindexentry{反事实鲁棒性}{Counterfactual Robustness}{\bookindexpage{section.7.5}{294}, \bookindexpage{section.7.5}{305–306}}
\bookindexentry{反事实支持信息}{Counterfactual Supporting Information}{\bookindexpage{section.7.5}{306}}
\bookindexentry{}{CRITIC}{\bookindexpage{section.6.1}{228}}
\bookindexentry{交叉注意力}{Cross-Attention}{\bookindexpage{section.9.3}{382}}
\bookindexentry{跨流架构}{Cross-Stream Architecture}{\bookindexpage{section.9.3}{380}}
\bookindexentry{CrowS-Pairs 分数}{CrowS-Pairs Score}{\bookindexpage{section.6.2}{235}}
\bookindexgroup{D}
\bookindexentry{}{DALL-E 2}{\bookindexpage{section.9.3}{377}}
\bookindexentry{数据治理}{Data Governance}{\bookindexpage{section.8.8}{361}}
\bookindexentry{数据并行}{Data Parallelism}{\bookindexpage{section.2.5}{60}}
\bookindexentry{数据安全}{Data Security}{\bookindexpage{section.8.2}{316}, \bookindexpage{section.8.5}{341–342}}
\bookindexentry{数据版本管理}{Data Versioning}{\bookindexpage{section.8.8}{356}}
\bookindexentry{解码器}{Decoder}{\bookindexpage{section.2.3}{34}}
\bookindexentry{}{DeepMind}{\bookindexpage{section.5.4}{194}}
\bookindexentry{}{DeepSpeed}{\bookindexpage{section.2.5}{60}}
\bookindexentry{人口群体表征}{Demographic Representation}{\bookindexpage{section.6.2}{236}}
\bookindexentry{去噪自编码器}{Denoising Autoencoder}{\bookindexpage{section.2.5}{55}}
\bookindexentry{深度优先搜索}{Depth-first Search}{\bookindexpage{section.4.6}{168}}
\bookindexentry{差分隐私}{Differential Privacy}{\bookindexpage{section.6.4}{260}}
\bookindexentry{}{Dify.Ai}{\bookindexpage{section.8.6}{345}}
\bookindexentry{}{DisCo}{\bookindexpage{section.6.2}{234}}
\bookindexentry{离散提示}{Discrete Prompt}{\bookindexpage{section.3.3}{107}}
\bookindexentry{}{DistilBERT}{\bookindexpage{section.2.6}{64}}
\bookindexentry{}{DistilGPT-2}{\bookindexpage{section.4.6}{169}}
\bookindexentry{}{DLite}{\bookindexpage{section.5.6}{214}}
\bookindexentry{领域适配}{Domain Adaptation}{\bookindexpage{section.4.2}{142}–\bookindexpage{section.4.3}{144}, \bookindexpage{section.6.3}{251}}
\bookindexentry{领域专用语言模型}{Domain-Specific Language Models}{\bookindexpage{section.10.3}{430}–\bookindexpage{section.10.4}{431}}
\bookindexentry{点积}{Dot Product}{\bookindexpage{section.2.2}{32}}
\bookindexentry{直接偏好优化}{DPO, Direct Preference Optimization}{\bookindexpage{section.5.5}{210–212}}
\bookindexentry{}{DSPy}{\bookindexpage{section.8.6}{345}}
\bookindexentry{动态记忆网络}{Dynamic Memory Network}{\bookindexpage{section.3.1}{84}}
\bookindexgroup{E}
\bookindexentry{早期拼接}{Early Concatenation}{\bookindexpage{section.9.3}{382}}
\bookindexentry{早期求和}{Early Summation}{\bookindexpage{section.9.3}{382}}
\bookindexentry{}{ELECTRA}{\bookindexpage{section.2.6}{64}}
\bookindexentry{}{Elo}{\bookindexpage{section.5.5}{209}}
\bookindexentry{嵌入模型}{Embedding Models}{\bookindexpage{section.7.3}{280}, \bookindexpage{section.7.4}{283}}
\bookindexentry{嵌入}{Embeddings}{\bookindexpage{section.6.2}{233}, \bookindexpage{section.8.6}{346}}
\bookindexentry{涌现能力}{Emergent Abilities}{\bookindexpage{section.1.4}{12}–\bookindexpage{section.1.5}{13}}
\bookindexentry{编码器}{Encoder}{\bookindexpage{section.2.3}{33}}
\bookindexentry{编码器架构}{Encoder Architecture}{\bookindexpage{section.2.5}{57}}
\bookindexentry{编码器—解码器架构}{Encoder Decoder Architecture}{\bookindexpage{section.1.5}{16}, \bookindexpage{section.2.1}{29}, \bookindexpage{section.2.5}{55}}
\bookindexentry{编码器—解码器注意力}{Encoder-Decoder Attention}{\bookindexpage{section.2.3}{38}}
\bookindexentry{纯编码器架构}{Encoder-Only Architecture}{\bookindexpage{section.1.5}{15}, \bookindexpage{section.2.1}{30}}
\bookindexentry{纯编码器架构}{Encoder-only Architecture}{\bookindexpage{section.2.1}{30}}
\bookindexentry{终端用户提示}{End-User Prompting}{\bookindexpage{section.1.5}{20}}
\bookindexentry{实体标注}{Entity Tagging}{\bookindexpage{section.3.2}{96}}
\bookindexentry{纠错率}{Error Correction Rate}{\bookindexpage{section.7.5}{306}}
\bookindexentry{错误检出率}{Error Detection Rate}{\bookindexpage{section.7.5}{306}}
\bookindexentry{}{Eva-CLIP-ViT}{\bookindexpage{section.9.3}{378}}
\bookindexentry{精确匹配}{Exact Match}{\bookindexpage{section.7.5}{302}, \bookindexpage{section.7.5}{304}}
\bookindexentry{期望最大化算法}{Expectation-Maximization Algorithm}{\bookindexpage{section.1.3}{8}}
\bookindexentry{期望最大毒性}{Expected Maximum Toxicity}{\bookindexpage{section.6.3}{249}}
\bookindexentry{实验跟踪}{Experiment Tracking}{\bookindexpage{section.8.8}{356}}
\bookindexentry{梯度爆炸}{Exploding Gradients}{\bookindexpage{section.2.1}{31}}
\bookindexgroup{F}
\bookindexentry{F1 分数}{F1 Score}{\bookindexpage{section.8.3}{326}}
\bookindexentry{事实检索}{Fact Retrieval}{\bookindexpage{section.3.3}{107}}
\bookindexentry{事实验证}{Fact Verification}{\bookindexpage{section.6.1}{222}}
\bookindexentry{}{FacTool}{\bookindexpage{section.6.1}{228}}
\bookindexentry{}{FActScore}{\bookindexpage{section.6.1}{221}, \bookindexpage{section.6.1}{225}}
\bookindexentry{事实一致性}{Factual Consistency}{\bookindexpage{section.6.1}{222}}
\bookindexentry{}{FactualityPrompt}{\bookindexpage{section.6.1}{223}}
\bookindexentry{公平性}{Fairness}{\bookindexpage{section.6.2}{233}}
\bookindexentry{}{FAISS}{\bookindexpage{section.8.6}{346}}
\bookindexentry{忠实度}{Faithfulness}{\bookindexpage{section.7.4}{284}}
\bookindexentry{前馈神经网络}{Feed-Forward Neural Network}{\bookindexpage{section.2.3}{36}, \bookindexpage{section.2.3}{42}}
\bookindexentry{反馈模型}{Feedback Model}{\bookindexpage{section.5.5}{208}}
\bookindexentry{少样本}{Few-Shot}{\bookindexpage{section.3.1}{89}, \bookindexpage{section.7.4}{285}}
\bookindexentry{少样本提示}{Few-Shot Prompts}{\bookindexpage{section.4.5}{164}}
\bookindexentry{金融服务}{Financial Services}{\bookindexpage{section.8.5}{342}}
\bookindexentry{}{FinBERT}{\bookindexpage{section.10.4}{431}}
\bookindexentry{微调}{Fine-tuning}{\bookindexpage{section.1.5}{19}, \bookindexpage{section.7.5}{305}, \bookindexpage{section.8.5}{340}}
\bookindexentry{}{Flamingo}{\bookindexpage{section.9.5}{401–405}}
\bookindexentry{}{Flan-PaLM}{\bookindexpage{section.10.4}{431}}
\bookindexentry{}{FLARE}{\bookindexpage{section.10.4}{434}}
\bookindexentry{}{Flowise}{\bookindexpage{section.8.6}{344}}
\bookindexentry{帧顺序建模}{Frame Order Modeling}{\bookindexpage{section.9.3}{390}}
\bookindexgroup{G}
\bookindexentry{}{G-Eval}{\bookindexpage{section.8.3}{323}, \bookindexpage{section.8.3}{328–329}}
\bookindexentry{}{Galactica}{\bookindexpage{section.10.4}{431}}
\bookindexentry{生成增强检索}{GAR, Generation-Augmented Retrieval}{\bookindexpage{section.7.4}{285}}
\bookindexentry{门控状态空间}{Gated State Space}{\bookindexpage{section.10.3}{428}}
\bookindexentry{}{GatorTron}{\bookindexpage{section.10.4}{431}}
\bookindexentry{GeGLU 激活}{GeGLU Activation}{\bookindexpage{section.2.3}{41}}
\bookindexentry{}{Gemma}{\bookindexpage{section.2.6}{73}, \bookindexpage{section.10.3}{429}}
\bookindexentry{生成式判别器}{Generative Discriminator}{\bookindexpage{section.6.3}{253}}
\bookindexentry{几何平均}{Geometric Mean}{\bookindexpage{section.8.3}{324}}
\bookindexentry{带音频反馈的手势控制}{Gesture-Based Control with Audio Feedback}{\bookindexpage{section.9.2}{376}}
\bookindexentry{Git 版本控制}{git (version control)}{\bookindexpage{section.8.8}{356}}
\bookindexentry{}{Github Copilot}{\bookindexpage{section.8.2}{318}}
\bookindexentry{}{Github Gists}{\bookindexpage{section.8.2}{320}}
\bookindexentry{}{GLUE}{\bookindexpage{section.2.6}{67}}
\bookindexentry{}{Google Vertex AI}{\bookindexpage{section.8.6}{346}}
\bookindexentry{}{Gopher}{\bookindexpage{section.2.3}{39}, \bookindexpage{section.2.6}{72}}
\bookindexentry{}{GPipe}{\bookindexpage{section.10.2}{426}}
\bookindexentry{}{GPT-1}{\bookindexpage{section.2.6}{68}, \bookindexpage{section.2.6}{69}}
\bookindexentry{}{GPT-2}{\bookindexpage{section.1.4}{11}, \bookindexpage{section.2.6}{68–70}, \bookindexpage{section.2.7}{77}, \bookindexpage{section.3.3}{106}, \bookindexpage{section.3.3}{110}, \bookindexpage{section.4.4}{157}, \bookindexpage{section.5.4}{196}, \bookindexpage{section.6.3}{251}, \bookindexpage{section.6.4}{256}, \bookindexpage{section.9.5}{401}}
\bookindexentry{}{GPT-3}{\bookindexpage{section.1.4}{11}, \bookindexpage{section.2.5}{58}, \bookindexpage{section.2.6}{68–70}, \bookindexpage{section.3.1}{84}, \bookindexpage{section.4.2}{141}, \bookindexpage{section.4.2}{142}, \bookindexpage{section.4.3}{155}, \bookindexpage{section.4.4}{158}, \bookindexpage{section.4.5}{166}, \bookindexpage{section.5.1}{178}, \bookindexpage{section.6.1}{222}}
\bookindexentry{}{GPT-3.5}{\bookindexpage{section.2.6}{68}, \bookindexpage{section.2.6}{73}, \bookindexpage{section.7.3}{280}}
\bookindexentry{}{GPT-4}{\bookindexpage{section.4.6}{169}, \bookindexpage{section.6.1}{220}, \bookindexpage{section.6.4}{259}, \bookindexpage{section.7.3}{280}, \bookindexpage{section.7.4}{283}, \bookindexpage{section.7.5}{299}, \bookindexpage{section.8.6}{347}, \bookindexpage{section.8.9}{362}, \bookindexpage{section.9.4}{399}}
\bookindexentry{}{GPT-4V}{\bookindexpage{section.9.4}{399}}
\bookindexentry{}{GPT-Neo-6}{\bookindexpage{section.6.4}{257}}
\bookindexentry{}{GPT4All}{\bookindexpage{section.6.5}{262}}
\bookindexentry{梯度检查点}{Gradient Checkpointing}{\bookindexpage{section.10.2}{426}}
\bookindexentry{梯度下降}{Gradient Descent}{\bookindexpage{section.5.5}{210}}
\bookindexentry{贪心搜索}{Greedy Search}{\bookindexpage{section.2.5}{58–59}, \bookindexpage{section.5.5}{210}}
\bookindexentry{真实答案}{Ground-truth Answers}{\bookindexpage{section.7.5}{301}}
\bookindexentry{有据性}{Groundedness}{\bookindexpage{section.5.4}{189}, \bookindexpage{section.5.4}{190}, \bookindexpage{section.5.4}{193}, \bookindexpage{section.7.5}{298}}
\bookindexentry{}{GSM8K}{\bookindexpage{section.4.5}{166}}
\bookindexgroup{H}
\bookindexentry{幻觉}{Hallucination}{\bookindexpage{section.4.5}{167}, \bookindexpage{section.6.1}{219}–\bookindexpage{section.6.2}{229}, \bookindexpage{section.7.1}{275}, \bookindexpage{section.7.4}{285}, \bookindexpage{section.7.5}{297}, \bookindexpage{section.9.4}{399}}
\bookindexentry{}{HaluEval}{\bookindexpage{section.6.1}{225}}
\bookindexentry{无害性}{Harmlessness}{\bookindexpage{section.5.1}{178}, \bookindexpage{section.5.1}{182–183}, \bookindexpage{section.5.5}{206}}
\bookindexentry{基于任务头的调优}{Head-based Tuning}{\bookindexpage{section.3.2}{91}}
\bookindexentry{医疗}{Healthcare}{\bookindexpage{section.8.5}{342}}
\bookindexentry{有帮助性}{Helpfulness}{\bookindexpage{section.5.1}{178–180}, \bookindexpage{section.5.5}{209}}
\bookindexentry{}{HERO}{\bookindexpage{section.9.3}{389}}
\bookindexentry{分层注意力}{Hierarchical Attention}{\bookindexpage{section.9.3}{382}}
\bookindexentry{命中率}{Hit Rate}{\bookindexpage{section.7.5}{296}}
\bookindexentry{诚实性指标}{Honest Metric}{\bookindexpage{section.6.2}{237}}
\bookindexentry{诚实性}{Honesty}{\bookindexpage{section.5.1}{178}, \bookindexpage{section.5.1}{181–182}}
\bookindexentry{}{HuBERT}{\bookindexpage{section.9.3}{379}}
\bookindexentry{}{HuggingFace}{\bookindexpage{section.3.1}{87}, \bookindexpage{section.4.3}{147}, \bookindexpage{section.8.2}{320}, \bookindexpage{section.8.6}{344}, \bookindexpage{section.10.3}{430}}
\bookindexentry{人工评估}{Human Evaluation}{\bookindexpage{section.6.1}{221–222}, \bookindexpage{section.8.3}{322}}
\bookindexentry{人类偏好}{Human Preferences}{\bookindexpage{section.5.5}{212}}
\bookindexentry{人工判断}{Human-based Judgement}{\bookindexpage{section.7.5}{298}}
\bookindexentry{人工生成内容}{Human-Generated Content}{\bookindexpage{section.8.2}{318}}
\bookindexentry{人在回路中的流程}{Human-in-the-Loop Processes}{\bookindexpage{section.8.8}{359}}
\bookindexentry{HurtLex 语料库}{HurtLex corpus}{\bookindexpage{section.6.2}{237}}
\bookindexentry{假设文档编码}{HyDE, Hypothetical Document Encoding}{\bookindexpage{section.7.4}{286}}
\bookindexgroup{I}
\bookindexentry{上下文学习}{ICL, In-Context Learning}{\bookindexpage{section.1.5}{20–21}, \bookindexpage{section.3.3}{110}, \bookindexpage{section.3.5}{118–119}, \bookindexpage{section.8.5}{341}}
\bookindexentry{图像和视频描述生成}{Image and Video Captioning}{\bookindexpage{section.9.2}{376}}
\bookindexentry{图像分类}{Image Classification}{\bookindexpage{section.9.3}{377}}
\bookindexentry{图像到文本的生成}{Image-Text Generation}{\bookindexpage{section.9.3}{388–389}}
\bookindexentry{图文匹配}{Image-Text Matching}{\bookindexpage{section.9.3}{388}}
\bookindexentry{图像—文本模态}{Image-Text Modality}{\bookindexpage{section.9.3}{384}}
\bookindexentry{图像与文本检索}{Image-Text Retrieval}{\bookindexpage{section.9.2}{376}}
\bookindexentry{}{ImageBind}{\bookindexpage{section.9.3}{377}}
\bookindexentry{训练中缓解}{In-training Mitigation}{\bookindexpage{section.6.2}{242–244}}
\bookindexentry{索引存储}{Index Storage}{\bookindexpage{section.7.3}{280}}
\bookindexentry{索引构建}{Indexing}{\bookindexpage{section.7.2}{276}}
\bookindexentry{推理延迟}{Inference Latency}{\bookindexpage{section.8.5}{337–339}}
\bookindexentry{}{Infini-attention}{\bookindexpage{section.10.2}{425}}
\bookindexentry{信息整合}{Information Integration}{\bookindexpage{section.7.5}{294}, \bookindexpage{section.7.5}{303–304}}
\bookindexentry{信息检索}{Information Retrieval}{\bookindexpage{section.8.2}{318}}
\bookindexentry{输入投影器}{Input Projector}{\bookindexpage{section.9.3}{379–383}}
\bookindexentry{}{Instruct-FinGPT}{\bookindexpage{section.4.2}{143}}
\bookindexentry{}{InstructGPT}{\bookindexpage{section.1.4}{11}, \bookindexpage{section.2.6}{68}, \bookindexpage{section.5.4}{200}, \bookindexpage{section.5.6}{213}}
\bookindexentry{}{InstructGPT}{\bookindexpage{section.4.2}{141–142}}
\bookindexentry{指令调优}{Instruction Tuning}{\bookindexpage{section.1.5}{18}, \bookindexpage{section.4.2}{136}–\bookindexpage{section.4.3}{144}, \bookindexpage{section.8.8}{358}}
\bookindexentry{}{INT8}{\bookindexpage{section.4.4}{158}, \bookindexpage{section.4.4}{159}}
\bookindexentry{}{InterBERT}{\bookindexpage{section.9.3}{382}}
\bookindexentry{内在维度}{Intrinsic Dimensions}{\bookindexpage{section.4.3}{150}, \bookindexpage{section.4.3}{153}, \bookindexpage{section.4.3}{155}}
\bookindexentry{迭代检索}{Iterative Retrieval}{\bookindexpage{section.7.4}{289}}
\bookindexgroup{J}
\bookindexentry{越狱}{Jail-Breaking}{\bookindexpage{section.5.2}{184}}
\bookindexentry{越狱}{Jailbreak}{\bookindexpage{section.8.5}{342}}
\bookindexentry{}{JavaScript}{\bookindexpage{section.8.2}{319}, \bookindexpage{section.8.6}{344}}
\bookindexentry{}{JetBrains}{\bookindexpage{section.8.2}{320}}
\bookindexgroup{K}
\bookindexentry{知识 F1}{K-F1; Knowledge-F1}{\bookindexpage{section.7.5}{298}}
\bookindexentry{知识精确率}{K-Precision, Knowledge-Precision}{\bookindexpage{section.7.5}{298}}
\bookindexentry{Kaiming 初始化}{Kaiming Initialization}{\bookindexpage{section.4.4}{157}}
\bookindexentry{键值缓存}{Key-Value Caching}{\bookindexpage{section.8.7}{353}}
\bookindexentry{知识库}{Knowledge Base}{\bookindexpage{section.3.2}{95}, \bookindexpage{section.6.1}{228}, \bookindexpage{section.7.1}{275}, \bookindexpage{section.7.4}{281}, \bookindexpage{section.7.5}{305}}
\bookindexentry{知识蒸馏}{Knowledge Distillation}{\bookindexpage{section.10.2}{427}–\bookindexpage{section.10.3}{428}}
\bookindexentry{知识抽取}{Knowledge Extraction}{\bookindexpage{section.3.2}{94}}
\bookindexentry{知识继承预训练}{Knowledge Inherited Pre-training}{\bookindexpage{section.2.5}{50}}
\bookindexentry{}{KOSMOS-1}{\bookindexpage{section.9.3}{377}}
\bookindexentry{克罗内克适配器}{Kronecker Adapters}{\bookindexpage{section.4.3}{151}, \bookindexpage{section.4.3}{153–154}}
\bookindexentry{Kullback–Leibler 散度}{Kullback-Leibler Divergence}{\bookindexpage{section.5.4}{199}, \bookindexpage{section.5.5}{212}, \bookindexpage{section.6.2}{237}}
\bookindexgroup{L}
\bookindexentry{}{LaMDA}{\bookindexpage{section.5.4}{189–191}}
\bookindexentry{}{LAMM}{\bookindexpage{section.9.5}{400}}
\bookindexentry{}{LangChain}{\bookindexpage{section.7.4}{288}, \bookindexpage{section.8.6}{344}, \bookindexpage{section.10.4}{433}}
\bookindexentry{}{LangKit}{\bookindexpage{section.8.9}{363}}
\bookindexentry{语言建模}{Language Modeling}{\bookindexpage{section.1.3}{7–8}}
\bookindexentry{层归一化}{Layer Normalization}{\bookindexpage{section.2.3}{37}}
\bookindexentry{层归一化}{LayerNorm}{\bookindexpage{section.2.3}{39}}
\bookindexentry{学习率}{Learning Rate}{\bookindexpage{section.2.5}{58}}
\bookindexentry{基于词语的指标}{Lexical-based Metrics}{\bookindexpage{section.7.5}{298}}
\bookindexentry{许可}{Licensing}{\bookindexpage{section.8.2}{316}, \bookindexpage{section.8.5}{341–342}}
\bookindexentry{}{LIME}{\bookindexpage{section.6.3}{247}}
\bookindexentry{}{Llama}{\bookindexpage{section.2.6}{72}, \bookindexpage{section.4.2}{143}–\bookindexpage{section.4.3}{144}}
\bookindexentry{}{Llama-2}{\bookindexpage{section.2.6}{73}, \bookindexpage{section.2.6}{75}, \bookindexpage{section.5.1}{179}, \bookindexpage{section.5.1}{180}, \bookindexpage{section.7.3}{280}, \bookindexpage{section.8.6}{347}}
\bookindexentry{}{Llama-3}{\bookindexpage{section.10.3}{429}}
\bookindexentry{}{LlamaIndex}{\bookindexpage{section.7.2}{276}, \bookindexpage{section.7.4}{283}, \bookindexpage{section.7.4}{288}, \bookindexpage{section.7.4}{291}, \bookindexpage{section.7.5}{293}, \bookindexpage{section.7.6}{307}, \bookindexpage{section.8.6}{344}, \bookindexpage{section.10.4}{433}}
\bookindexentry{大语言模型智能体}{LLM Agents}{\bookindexpage{section.10.4}{431–433}}
\bookindexentry{大语言模型评估}{LLM Evaluation}{\bookindexpage{section.8.3}{323}}
\bookindexentry{}{LLM.int8}{\bookindexpage{section.4.4}{160}}
\bookindexentry{}{LLMCritic}{\bookindexpage{section.7.5}{299}}
\bookindexentry{}{LLMDataHub}{\bookindexpage{section.10.3}{430}}
\bookindexentry{大语言模型运维}{LLMOps}{\bookindexpage{section.8.2}{316}, \bookindexpage{section.8.6}{345}, \bookindexpage{section.8.8}{354–361}}
\bookindexentry{}{LoftQ}{\bookindexpage{section.4.4}{161}–\bookindexpage{section.4.5}{162}}
\bookindexentry{对数概率偏见分数}{Log-Probability Bias Score}{\bookindexpage{section.6.2}{235}}
\bookindexentry{逻辑回归}{Logistic Regression}{\bookindexpage{section.1.3}{8}}
\bookindexentry{最长公共子序列}{Longest Common Subsequence}{\bookindexpage{section.8.3}{325}}
\bookindexentry{}{LongRoPE}{\bookindexpage{section.10.2}{425}}
\bookindexentry{低秩适配}{LoRA}{\bookindexpage{section.1.5}{19}, \bookindexpage{section.4.3}{145}, \bookindexpage{section.4.3}{151–153}, \bookindexpage{section.4.3}{155}, \bookindexpage{section.4.4}{161}, \bookindexpage{section.4.6}{168}}
\bookindexentry{}{LoRAHub}{\bookindexpage{section.4.3}{152}}
\bookindexgroup{M}
\bookindexentry{}{M3DBench}{\bookindexpage{section.9.4}{399}}
\bookindexentry{机器翻译}{Machine Translation}{\bookindexpage{section.2.5}{55}, \bookindexpage{section.4.3}{149}}
\bookindexentry{手工答案映射}{Manual Answer Mapping}{\bookindexpage{section.3.4}{115}}
\bookindexentry{手工提示}{Manual Prompt}{\bookindexpage{section.3.3}{101–103}, \bookindexpage{section.6.2}{242}}
\bookindexentry{马尔可夫决策过程}{Markov Decision Process}{\bookindexpage{section.5.2}{186}–\bookindexpage{section.5.3}{188}}
\bookindexentry{掩码语言建模}{Masked Language Modeling}{\bookindexpage{section.2.5}{51}, \bookindexpage{section.2.6}{62}, \bookindexpage{section.9.3}{387}}
\bookindexentry{掩码多头注意力}{Masked Multi-Headed Attention}{\bookindexpage{section.2.3}{38}}
\bookindexentry{掩码物体分类}{Masked Object Classification}{\bookindexpage{section.9.3}{387–388}}
\bookindexentry{}{Megatron-LM}{\bookindexpage{section.2.7}{76}}
\bookindexentry{元宇宙}{Metaverse}{\bookindexpage{section.9.3}{379}}
\bookindexentry{}{Milvus}{\bookindexpage{section.8.6}{346}}
\bookindexentry{}{Mistral 7B}{\bookindexpage{section.8.3}{322}, \bookindexpage{section.10.3}{429}}
\bookindexentry{}{Mixtral}{\bookindexpage{section.2.6}{73}–\bookindexpage{section.2.7}{76}}
\bookindexentry{混合专家}{Mixture of Experts}{\bookindexpage{section.2.5}{57}, \bookindexpage{section.2.6}{73}}
\bookindexentry{最大似然估计}{MLE, Maximum Likelihood Estimation}{\bookindexpage{section.5.5}{212}}
\bookindexentry{机器学习运维}{MLOps}{\bookindexpage{section.8.8}{354}}
\bookindexentry{}{MMBench}{\bookindexpage{section.9.4}{399}}
\bookindexentry{}{MMCBench}{\bookindexpage{section.9.4}{399}}
\bookindexentry{}{MMHAL-BENCH}{\bookindexpage{section.9.4}{399}}
\bookindexentry{}{MMMU}{\bookindexpage{section.9.4}{398}}
\bookindexentry{}{MMVP}{\bookindexpage{section.9.4}{399}}
\bookindexentry{模态对齐}{Modality Alignment}{\bookindexpage{section.9.3}{379}}
\bookindexentry{模态编码器}{Modality Encoder}{\bookindexpage{section.9.3}{378–379}}
\bookindexentry{模态生成器}{Modality Generator}{\bookindexpage{section.9.4}{398}}
\bookindexentry{模态匹配损失}{Modality Matching Loss}{\bookindexpage{section.9.3}{386–387}}
\bookindexentry{模态切换指令调优}{Modality-Switching Instruction Tuning}{\bookindexpage{section.9.5}{410}}
\bookindexentry{模型治理}{Model Governance}{\bookindexpage{section.8.8}{361}}
\bookindexentry{基于模型的指标}{Model-based Metrics}{\bookindexpage{section.7.5}{298–300}}
\bookindexentry{模块化 RAG}{Modular RAG}{\bookindexpage{section.7.4}{281}}
\bookindexentry{混合专家}{MOE, Mixture of Experts}{\bookindexpage{section.2.6}{74}–\bookindexpage{section.2.7}{76}}
\bookindexentry{监控}{Monitoring}{\bookindexpage{section.8.8}{357–358}}
\bookindexentry{}{MoverScore}{\bookindexpage{section.8.3}{326–328}}
\bookindexentry{平均倒数排名}{MRR, Mean Reciprocal Rank}{\bookindexpage{section.7.5}{296}}
\bookindexentry{}{MS COCO}{\bookindexpage{section.9.5}{410}}
\bookindexentry{多头注意力}{Multi-Head Attention}{\bookindexpage{section.2.3}{36}}
\bookindexentry{多跳问题}{Multi-Hop Question}{\bookindexpage{section.7.4}{289}}
\bookindexentry{多提示推理}{Multi-Prompt Inference}{\bookindexpage{section.3.5}{116}–\bookindexpage{section.3.6}{120}}
\bookindexentry{多任务学习}{Multi-task Learning}{\bookindexpage{section.2.6}{65}}
\bookindexentry{多轮对话}{Multi-turn Conversation}{\bookindexpage{section.5.4}{190}}
\bookindexentry{多语言大语言模型}{Multilingual Large Language Models}{\bookindexpage{section.8.2}{321}}
\bookindexentry{多模态思维链}{Multimodal Chain-of-Thought}{\bookindexpage{section.9.3}{391}, \bookindexpage{section.9.3}{394–395}}
\bookindexentry{多模态上下文学习}{Multimodal In-Context Learning}{\bookindexpage{section.9.3}{392–394}}
\bookindexentry{多模态指令调优}{Multimodal Instruction Tuning}{\bookindexpage{section.9.3}{391–392}}
\bookindexentry{多模态大语言模型}{Multimodal Large Language Models}{\bookindexpage{section.8.2}{321}, \bookindexpage{section.9.1}{375}–\bookindexpage{section.9.6}{415}}
\bookindexentry{多模态 RLHF}{Multimodal RLHF}{\bookindexpage{section.9.3}{395–397}}
\bookindexentry{}{MultiNLI}{\bookindexpage{section.2.6}{68}}
\bookindexentry{多词元生成}{Multitoken Generation}{\bookindexpage{section.10.2}{426–427}}
\bookindexentry{}{MXNet}{\bookindexpage{section.10.2}{426}}
\bookindexgroup{N}
\bookindexentry{自然语言}{Natural Language}{\bookindexpage{section.1.2}{3}–\bookindexpage{section.1.3}{5}}
\bookindexentry{自然语言推断}{Natural Language Inference}{\bookindexpage{section.2.6}{61}, \bookindexpage{section.3.3}{107}, \bookindexpage{section.6.1}{222}, \bookindexpage{section.7.5}{299}}
\bookindexentry{自然语言理解}{Natural Language Understanding}{\bookindexpage{section.1.3}{5}, \bookindexpage{section.8.2}{319}}
\bookindexentry{归一化折损累积增益}{nDCG, Normalized Discounted Cumulative Gain}{\bookindexpage{section.7.5}{297}}
\bookindexentry{负例上下文}{Negative Contexts}{\bookindexpage{section.7.5}{302}}
\bookindexentry{无答案时拒答}{Negative Rejection}{\bookindexpage{section.7.5}{294}, \bookindexpage{section.7.5}{302–303}}
\bookindexentry{命名实体识别}{NER, Named-Entity Recognition}{\bookindexpage{section.1.4}{10}, \bookindexpage{section.1.5}{18}, \bookindexpage{section.7.5}{299}}
\bookindexentry{}{Neural-Chat-v3-1}{\bookindexpage{section.10.2}{424}}
\bookindexentry{《纽约时报》}{New York Times}{\bookindexpage{section.8.5}{342}}
\bookindexentry{下一句预测}{Next Sentence Prediction}{\bookindexpage{section.2.5}{54}, \bookindexpage{section.2.6}{62}}
\bookindexentry{}{NExT-GPT}{\bookindexpage{section.9.3}{377}, \bookindexpage{section.9.5}{401}, \bookindexpage{section.9.5}{408}–\bookindexpage{section.9.6}{411}}
\bookindexentry{自然语言处理}{NLP, Natural Language Processing}{\bookindexpage{section.1.3}{5–8}}
\bookindexentry{噪声鲁棒性}{Noise Robustness}{\bookindexpage{section.7.5}{294}, \bookindexpage{section.7.5}{301–302}}
\bookindexentry{}{Nomic AI}{\bookindexpage{section.6.5}{262}}
\bookindexentry{NVIDIA Triton 推理服务器}{NVIDIA’s Triton Inference Server}{\bookindexpage{section.8.7}{352}}
\bookindexgroup{O}
\bookindexentry{光学字符识别}{OCR}{\bookindexpage{section.9.4}{399}}
\bookindexentry{单样本}{One-Shot}{\bookindexpage{section.3.1}{89}}
\bookindexentry{本体知识}{Ontological Knowledge}{\bookindexpage{section.6.1}{220}}
\bookindexentry{}{OpaquePrompts}{\bookindexpage{section.6.4}{261}}
\bookindexentry{}{OpenAI}{\bookindexpage{section.2.6}{68}, \bookindexpage{section.4.2}{140}, \bookindexpage{section.5.4}{196}, \bookindexpage{section.6.4}{261}, \bookindexpage{section.7.2}{277}, \bookindexpage{section.7.6}{308}, \bookindexpage{section.8.2}{318}}
\bookindexentry{}{OpenHermes}{\bookindexpage{section.10.2}{424}}
\bookindexentry{}{OpenLLM}{\bookindexpage{section.10.3}{430}}
\bookindexentry{}{OPENWEBTEXTCORPUS}{\bookindexpage{section.6.3}{251}}
\bookindexentry{操作性条件作用}{Operant Conditioning}{\bookindexpage{section.5.2}{185}}
\bookindexentry{}{OPT-175B}{\bookindexpage{section.10.3}{430}}
\bookindexentry{输出投影器}{Output Projector}{\bookindexpage{section.9.3}{397}}
\bookindexentry{过拟合}{Overfitting}{\bookindexpage{section.6.2}{244}, \bookindexpage{section.6.4}{256}}
\bookindexgroup{P}
\bookindexentry{}{PaLM}{\bookindexpage{section.1.4}{11}}
\bookindexentry{并联适配器}{Parallel Adapters}{\bookindexpage{section.4.3}{146}, \bookindexpage{section.4.3}{149–150}, \bookindexpage{section.4.3}{154}}
\bookindexentry{并行性}{Parallelism}{\bookindexpage{section.8.5}{338}}
\bookindexentry{参数知识}{Parametric Knowledge}{\bookindexpage{section.3.2}{94}, \bookindexpage{section.6.1}{225}, \bookindexpage{section.6.1}{228}, \bookindexpage{section.7.5}{301}, \bookindexpage{section.7.5}{306}}
\bookindexentry{参数记忆}{Parametric Memory}{\bookindexpage{section.6.1}{228}}
\bookindexentry{pass@$k$}{Pass at k}{\bookindexpage{section.8.3}{329}}
\bookindexentry{基于提示的学习}{PBL, Prompt-based Learning}{\bookindexpage{section.4.2}{138}}
\bookindexentry{参数高效微调}{PEFT, Parameter-efficient fine-tuning}{\bookindexpage{section.1.5}{14}, \bookindexpage{section.1.5}{19–20}, \bookindexpage{section.4.3}{144}–\bookindexpage{section.4.4}{157}, \bookindexpage{section.9.3}{384}}
\bookindexentry{困惑度}{Perplexity}{\bookindexpage{section.6.4}{259}, \bookindexpage{section.7.4}{291}, \bookindexpage{section.8.3}{323}}
\bookindexentry{指定角色的提示}{Persona-Assigned Prompts}{\bookindexpage{section.6.3}{248}}
\bookindexentry{}{Phi-1.5}{\bookindexpage{section.6.1}{225}}
\bookindexentry{}{Phi-2}{\bookindexpage{section.10.3}{429}}
\bookindexentry{个人身份信息}{PII, Personally Identifiable Information}{\bookindexpage{section.6.4}{256}}
\bookindexentry{}{Pinecone}{\bookindexpage{section.8.6}{346}}
\bookindexentry{即插即用语言模型}{Plug and Play Language Model}{\bookindexpage{section.6.3}{253}}
\bookindexentry{位置编码}{Positional Encodings}{\bookindexpage{section.2.3}{35–36}}
\bookindexentry{位置跳跃式训练}{Positional Skip-wise Training}{\bookindexpage{section.10.2}{425}}
\bookindexentry{后处理缓解}{Post-processing Mitigation}{\bookindexpage{section.6.2}{246}–\bookindexpage{section.6.3}{247}}
\bookindexentry{训练后量化}{Post-training Quantization}{\bookindexpage{section.4.4}{159}}
\bookindexentry{}{Postgres}{\bookindexpage{section.8.6}{346}}
\bookindexentry{近端策略优化}{PPO, Proximal Policy Optimization}{\bookindexpage{section.5.4}{197}, \bookindexpage{section.6.1}{227}}
\bookindexentry{精确率}{Precision}{\bookindexpage{section.7.5}{295}, \bookindexpage{section.8.3}{324}, \bookindexpage{section.8.3}{326}}
\bookindexentry{}{Predibase}{\bookindexpage{section.8.8}{354}}
\bookindexentry{前缀解码器架构}{Prefix Decoder Architecture}{\bookindexpage{section.1.5}{16}, \bookindexpage{section.2.5}{56}}
\bookindexentry{前缀式提示}{Prefix Prompt}{\bookindexpage{section.3.3}{100–101}}
\bookindexentry{前缀调优}{Prefix Tuning}{\bookindexpage{section.3.3}{110}}
\bookindexentry{前缀语言建模}{PrefixLM}{\bookindexpage{section.9.3}{387}}
\bookindexentry{隐私}{Privacy}{\bookindexpage{section.5.1}{183}, \bookindexpage{section.6.4}{256}–\bookindexpage{section.6.5}{262}}
\bookindexentry{提示添加}{Prompt Addition}{\bookindexpage{section.3.2}{92–93}}
\bookindexentry{提示链}{Prompt Chaining}{\bookindexpage{section.4.5}{164–165}}
\bookindexentry{提示分解}{Prompt Decomposition}{\bookindexpage{section.3.5}{119}–\bookindexpage{section.3.6}{120}}
\bookindexentry{提示工程}{Prompt Engineering}{\bookindexpage{section.3.3}{97}–\bookindexpage{section.4.6}{168}, \bookindexpage{section.8.5}{341}, \bookindexpage{section.8.6}{347–349}}
\bookindexentry{提示集成}{Prompt Ensembling}{\bookindexpage{section.3.5}{116}}
\bookindexentry{}{Prompt Hub}{\bookindexpage{section.8.6}{349}}
\bookindexentry{提示挖掘}{Prompt Mining}{\bookindexpage{section.3.3}{103–104}}
\bookindexentry{提示改写}{Prompt Paraphrasing}{\bookindexpage{section.3.3}{104–105}}
\bookindexentry{提示版本管理}{Prompt Versioning}{\bookindexpage{section.8.8}{357}}
\bookindexentry{基于提示的微调}{Prompt-based fine-tuning}{\bookindexpage{section.3.4}{113}, \bookindexpage{section.4.2}{136}}
\bookindexentry{基于提示的学习}{Prompt-based Learning}{\bookindexpage{section.3.1}{88}–\bookindexpage{section.3.2}{92}}
\bookindexentry{}{Promptmetheus}{\bookindexpage{section.8.6}{348}}
\bookindexentry{}{ProPILE}{\bookindexpage{section.6.4}{257}}
\bookindexentry{伪对数似然}{Pseudo-log-likelihood}{\bookindexpage{section.6.2}{235}}
\bookindexentry{}{Pythia}{\bookindexpage{section.6.4}{256}}
\bookindexentry{}{pytorch}{\bookindexpage{section.8.6}{346}, \bookindexpage{section.10.2}{426}}
\bookindexgroup{Q}
\bookindexentry{}{QLoRA}{\bookindexpage{section.4.3}{153}, \bookindexpage{section.4.4}{161}, \bookindexpage{section.4.6}{168}}
\bookindexentry{量化}{Quantization}{\bookindexpage{section.4.4}{158}–\bookindexpage{section.4.5}{162}, \bookindexpage{section.8.5}{339}}
\bookindexentry{量化感知微调}{Quantization Aware Fine-tuning}{\bookindexpage{section.4.4}{160}–\bookindexpage{section.4.5}{162}}
\bookindexentry{量化感知训练}{Quantization-aware Training}{\bookindexpage{section.4.4}{158}}
\bookindexentry{查询改写}{Query Rewriting}{\bookindexpage{section.7.4}{284–285}}
\bookindexentry{查询到文档}{Query-to-Document}{\bookindexpage{section.7.4}{285}}
\bookindexentry{}{query2doc}{\bookindexpage{section.7.4}{286}}
\bookindexentry{问答}{Question Answering}{\bookindexpage{section.6.1}{222}}
\bookindexentry{问题自然语言推断}{Question NLI}{\bookindexpage{section.2.6}{68}}
\bookindexentry{问答对}{Question-Answer Pairs}{\bookindexpage{section.7.5}{303}}
\bookindexentry{问题—支持信息—答案三元组}{Question-Supporting Information-Answer Triplets}{\bookindexpage{section.7.5}{304}}
\bookindexgroup{R}
\bookindexentry{RAG 评估}{RAG Evaluation}{\bookindexpage{section.7.5}{293}–\bookindexpage{section.7.6}{307}}
\bookindexentry{检索增强生成}{RAG, Retrieval Augmented Generation}{\bookindexpage{section.1.5}{23}, \bookindexpage{section.7.2}{276}–\bookindexpage{section.7.6}{311}}
\bookindexentry{}{RAG-Fusion}{\bookindexpage{section.7.4}{285}}
\bookindexentry{}{Ragas}{\bookindexpage{section.7.5}{300}}
\bookindexentry{随机搜索}{Random Search}{\bookindexpage{section.2.5}{59–60}}
\bookindexentry{随机词元替换}{Random Token Substitution}{\bookindexpage{section.2.5}{52}}
\bookindexentry{与排名无关的指标}{Rank-Agnostic Metrics}{\bookindexpage{section.7.5}{294}}
\bookindexentry{考虑排名的指标}{Rank-Aware Metrics}{\bookindexpage{section.7.5}{294}, \bookindexpage{section.7.5}{297}}
\bookindexentry{}{RealToxicityPrompts}{\bookindexpage{section.5.4}{203}, \bookindexpage{section.6.2}{236}, \bookindexpage{section.6.2}{238}}
\bookindexentry{召回率}{Recall}{\bookindexpage{section.7.5}{294}, \bookindexpage{section.7.5}{295}, \bookindexpage{section.8.3}{326}}
\bookindexentry{倒数排名融合}{Reciprocal Rank Fusion}{\bookindexpage{section.7.4}{285}}
\bookindexentry{红队测试}{Red Teaming}{\bookindexpage{section.5.5}{206}}
\bookindexentry{}{Reddit}{\bookindexpage{section.5.4}{197}}
\bookindexentry{}{RedditBias}{\bookindexpage{section.6.2}{238}}
\bookindexentry{}{Reflexion}{\bookindexpage{section.10.4}{432}}
\bookindexentry{强化学习}{Reinforcement Learning}{\bookindexpage{section.5.2}{184}–\bookindexpage{section.5.3}{188}, \bookindexpage{section.5.4}{202}}
\bookindexentry{ReLU 激活}{ReLU Activation}{\bookindexpage{section.2.3}{37}}
\bookindexentry{重参数化}{Reparameterization}{\bookindexpage{section.4.3}{146}, \bookindexpage{section.4.3}{150}–\bookindexpage{section.4.4}{157}}
\bookindexentry{替换词元检测}{Replaced Token Detection}{\bookindexpage{section.2.5}{51}}
\bookindexentry{表征性伤害}{Representational Harm}{\bookindexpage{section.6.2}{230}}
\bookindexentry{奖励模型}{Reward Model}{\bookindexpage{section.5.3}{188}, \bookindexpage{section.5.4}{197–199}, \bookindexpage{section.5.4}{201}, \bookindexpage{section.6.1}{227}}
\bookindexentry{重加权}{Reweighting}{\bookindexpage{section.6.2}{240–241}}
\bookindexentry{风险管理}{Risk Management}{\bookindexpage{section.8.8}{359–360}}
\bookindexentry{基于人类反馈的强化学习}{RLHF, Reinforcement Learning from Human Feedback}{\bookindexpage{section.5.4}{189–204}}
\bookindexentry{循环神经网络}{RNNs, Recurrent Neural Networks}{\bookindexpage{section.2.1}{31}}
\bookindexentry{}{RoBERTa}{\bookindexpage{section.2.6}{64}, \bookindexpage{section.4.3}{155}, \bookindexpage{section.4.4}{157}}
\bookindexentry{}{ROUGE}{\bookindexpage{section.8.3}{323}, \bookindexpage{section.8.3}{325}}
\bookindexgroup{S}
\bookindexentry{安全性}{Safety}{\bookindexpage{section.5.4}{190}}
\bookindexentry{安全风险}{Safety Risks}{\bookindexpage{section.8.5}{341}}
\bookindexentry{}{SaulLM-7B}{\bookindexpage{section.10.4}{431}}
\bookindexentry{规模}{Scale}{\bookindexpage{section.1.4}{12}}
\bookindexentry{}{SciGraphQA}{\bookindexpage{section.9.5}{400}}
\bookindexentry{分段}{Segmentation}{\bookindexpage{section.7.4}{281}}
\bookindexentry{自注意力}{Self-Attention}{\bookindexpage{section.2.3}{33}}
\bookindexentry{自一致性}{Self-Consistency}{\bookindexpage{section.4.5}{167}}
\bookindexentry{自我去偏}{Self-Debiasing}{\bookindexpage{section.6.3}{254}}
\bookindexentry{自监督对比学习}{Self-Supervised Contrastive Learning}{\bookindexpage{section.9.3}{380}}
\bookindexentry{自监督语言建模}{Self-Supervised Language Modeling}{\bookindexpage{section.2.4}{43}}
\bookindexentry{自监督学习}{Self-Supervised Learning}{\bookindexpage{section.2.5}{50}}
\bookindexentry{}{SELFCHECKGPT}{\bookindexpage{section.6.2}{229}}
\bookindexentry{语义相似性}{Semantic Similarity}{\bookindexpage{section.7.4}{283}, \bookindexpage{section.7.4}{286}, \bookindexpage{section.7.4}{289}, \bookindexpage{section.7.5}{300}, \bookindexpage{section.8.3}{323}}
\bookindexentry{}{SemCor}{\bookindexpage{section.6.2}{232}}
\bookindexentry{合理性、具体性与趣味性}{Sensibleness, Specificity, and Interestingness (SSI)}{\bookindexpage{section.5.4}{189–191}}
\bookindexentry{句段边界目标}{Sentence Boundary Objective}{\bookindexpage{section.2.5}{53–54}}
\bookindexentry{句子编码器关联检验}{Sentence Encoder Association Test}{\bookindexpage{section.6.2}{234}}
\bookindexentry{句序预测}{Sentence Order Prediction}{\bookindexpage{section.2.5}{54}}
\bookindexentry{}{SentencePiece}{\bookindexpage{section.5.4}{190}}
\bookindexentry{情感分析}{Sentiment Analysis}{\bookindexpage{section.3.2}{94}}
\bookindexentry{情感分类}{Sentiment Classification}{\bookindexpage{section.3.2}{92}}
\bookindexentry{序列转换}{Sequence Transduction}{\bookindexpage{section.9.3}{379}}
\bookindexentry{串联适配器}{Series Adapters}{\bookindexpage{section.4.3}{146}, \bookindexpage{section.4.3}{148–149}}
\bookindexentry{}{SHAP}{\bookindexpage{section.6.3}{247}}
\bookindexentry{打乱词元检测}{Shuffled Token Detection}{\bookindexpage{section.2.5}{52}}
\bookindexentry{单流架构}{Single-Stream Architecture}{\bookindexpage{section.9.3}{380}}
\bookindexentry{小型语言模型}{Small Language Models}{\bookindexpage{section.10.3}{428}}
\bookindexentry{}{SNLI}{\bookindexpage{section.2.6}{68}}
\bookindexentry{社会偏见}{Social Bias}{\bookindexpage{section.6.2}{229}}
\bookindexentry{社会群体替换}{Social Group Substitutions}{\bookindexpage{section.6.2}{236}}
\bookindexentry{社会信号处理}{Social Signal Processing}{\bookindexpage{section.9.3}{377}}
\bookindexentry{}{SpanBERT}{\bookindexpage{section.2.6}{64}}
\bookindexentry{}{Sparkles}{\bookindexpage{section.9.5}{400}}
\bookindexentry{特殊词元}{Special Tokens}{\bookindexpage{section.2.6}{62}}
\bookindexentry{语音识别}{Speech Recognition}{\bookindexpage{section.9.3}{377}}
\bookindexentry{}{SQUAD}{\bookindexpage{section.2.6}{67}}
\bookindexentry{}{SqueezeLLM}{\bookindexpage{section.4.4}{159}, \bookindexpage{section.4.4}{160}}
\bookindexentry{}{Stable Diffusion}{\bookindexpage{section.9.6}{411}}
\bookindexentry{状态空间模型}{State Space Models}{\bookindexpage{section.10.3}{428}}
\bookindexentry{统计过滤}{Statistical Filtering}{\bookindexpage{section.2.4}{46}}
\bookindexentry{}{StereoSet}{\bookindexpage{section.6.2}{236}}
\bookindexentry{刻板印象关联}{Stereotypical Associations}{\bookindexpage{section.6.2}{237}}
\bookindexentry{存储成本}{Storage Costs}{\bookindexpage{section.8.5}{339}}
\bookindexentry{STS 基准}{STS Benchmark}{\bookindexpage{section.7.3}{280}}
\bookindexentry{摘要生成}{Summarization}{\bookindexpage{section.3.3}{110}, \bookindexpage{section.4.3}{151}, \bookindexpage{section.6.1}{222}}
\bookindexentry{}{SuperGLUE}{\bookindexpage{section.2.6}{67}, \bookindexpage{section.3.3}{112}}
\bookindexentry{}{Superlinked}{\bookindexpage{section.7.3}{280}}
\bookindexentry{监督微调}{Supervised Fine-tuning}{\bookindexpage{section.5.6}{214}}
\bookindexentry{监督学习}{Supervised Learning}{\bookindexpage{section.3.1}{85–86}}
\bookindexentry{支持向量机}{Support Vector Machines}{\bookindexpage{section.1.3}{8}}
\bookindexentry{替换式语言建模}{Swapped Language Modeling}{\bookindexpage{section.2.5}{52–53}}
\bookindexentry{SwiGLU 激活}{SwiGLU Activation}{\bookindexpage{section.2.3}{41}}
\bookindexentry{迎合}{Sycophancy}{\bookindexpage{section.6.1}{221}}
\bookindexentry{合成数据}{Synthetic Data}{\bookindexpage{section.10.2}{424}}
\bookindexgroup{T}
\bookindexentry{}{T5}{\bookindexpage{section.1.4}{11}, \bookindexpage{section.2.6}{64–68}, \bookindexpage{section.3.3}{108}, \bookindexpage{section.7.4}{285}}
\bookindexentry{任务适应性预训练}{Task Adaptive Pre-training}{\bookindexpage{section.2.5}{49}}
\bookindexentry{任务无关指标}{Task-Agnostic Metrics}{\bookindexpage{section.8.3}{322}}
\bookindexentry{任务专用指标}{Task-Specific Metrics}{\bookindexpage{section.8.3}{322}}
\bookindexentry{}{Teknium1}{\bookindexpage{section.10.2}{424}}
\bookindexentry{温度}{Temperature}{\bookindexpage{section.2.5}{59}, \bookindexpage{section.8.6}{347}}
\bookindexentry{张量}{Tensors}{\bookindexpage{section.2.6}{61}}
\bookindexentry{文本分块}{Text Chunking}{\bookindexpage{section.7.3}{279}}
\bookindexentry{文本分类}{Text Classification}{\bookindexpage{section.3.2}{96}, \bookindexpage{section.3.3}{108}}
\bookindexentry{文本生成}{Text Generation}{\bookindexpage{section.3.2}{96}}
\bookindexentry{}{TimeIT}{\bookindexpage{section.9.4}{399}}
\bookindexentry{}{TinyTransformer}{\bookindexpage{section.9.5}{410}}
\bookindexentry{词元切分}{Tokenization}{\bookindexpage{section.2.4}{47}}
\bookindexentry{Top-$k$ 采样}{Top-k Sampling}{\bookindexpage{section.2.5}{60}, \bookindexpage{section.5.5}{210}}
\bookindexentry{总拥有成本}{Total Cost of Ownership}{\bookindexpage{section.8.5}{333–335}, \bookindexpage{section.8.5}{340}}
\bookindexentry{毒性}{Toxicity}{\bookindexpage{section.6.3}{247–255}}
\bookindexentry{毒性概率}{Toxicity Probability}{\bookindexpage{section.6.3}{250}}
\bookindexentry{}{Transformers}{\bookindexpage{section.2.3}{33}–\bookindexpage{section.2.4}{43}}
\bookindexentry{翻译语言建模}{Translation Language Modeling}{\bookindexpage{section.2.5}{53}}
\bookindexentry{思维树}{Tree-of-Thoughts}{\bookindexpage{section.4.5}{167}–\bookindexpage{section.4.6}{168}}
\bookindexentry{触发词元}{Trigger Tokens}{\bookindexpage{section.3.3}{106}}
\bookindexentry{}{TruLens}{\bookindexpage{section.7.5}{300}, \bookindexpage{section.8.7}{350}, \bookindexpage{section.8.9}{363}}
\bookindexentry{信赖域策略优化}{Trust Region Policy Optimization}{\bookindexpage{section.5.4}{195}}
\bookindexentry{}{TrustGPT}{\bookindexpage{section.6.2}{238}}
\bookindexentry{真实性}{Truthfulness}{\bookindexpage{section.6.1}{228}}
\bookindexentry{}{TruthfulQA}{\bookindexpage{section.5.1}{181}, \bookindexpage{section.5.4}{203}, \bookindexpage{section.6.1}{221}, \bookindexpage{section.6.1}{225}, \bookindexpage{section.6.1}{226}}
\bookindexgroup{U}
\bookindexentry{}{ULIP}{\bookindexpage{section.9.3}{379}}
\bookindexentry{}{ULIP-2}{\bookindexpage{section.9.3}{379}}
\bookindexentry{}{ULMFiT}{\bookindexpage{section.2.4}{43}}
\bookindexentry{}{UltraFeedback}{\bookindexpage{section.10.2}{424}}
\bookindexgroup{V}
\bookindexentry{梯度消失}{Vanishing Gradients}{\bookindexpage{section.2.1}{31}}
\bookindexentry{向量数据库}{Vector Database}{\bookindexpage{section.7.5}{306}, \bookindexpage{section.8.6}{346–347}}
\bookindexentry{向量索引}{Vector Index}{\bookindexpage{section.7.4}{287}}
\bookindexentry{向量搜索}{Vector Search}{\bookindexpage{section.8.6}{347}}
\bookindexentry{}{VeRA}{\bookindexpage{section.4.3}{151}, \bookindexpage{section.4.3}{155}–\bookindexpage{section.4.4}{157}}
\bookindexentry{版本控制}{Version Control}{\bookindexpage{section.8.8}{356–357}}
\bookindexentry{}{Vicuna}{\bookindexpage{section.10.2}{424}, \bookindexpage{section.10.3}{429}}
\bookindexentry{}{Video-Bench}{\bookindexpage{section.9.4}{399}}
\bookindexentry{}{Video-LLaMA}{\bookindexpage{section.9.5}{401}, \bookindexpage{section.9.5}{405–407}}
\bookindexentry{视频与字幕匹配}{Video-Subtitle Matching}{\bookindexpage{section.9.3}{389–390}}
\bookindexentry{视频—文本模态}{Video-Text Modality}{\bookindexpage{section.9.3}{384}}
\bookindexentry{视频与文本检索}{Video-Text Retrieval}{\bookindexpage{section.9.2}{376}}
\bookindexentry{}{VilBERT}{\bookindexpage{section.9.3}{382}}
\bookindexentry{}{Vim}{\bookindexpage{section.8.2}{320}}
\bookindexentry{}{ViP-Bench}{\bookindexpage{section.9.4}{399}}
\bookindexentry{Visual Studio Code}{Visual Basic Code}{\bookindexpage{section.8.2}{320}}
\bookindexentry{视觉问答}{Visual Question Answering}{\bookindexpage{section.9.2}{376}}
\bookindexentry{视觉 Transformer}{ViT, Vision Transformer}{\bookindexpage{section.9.3}{378}}
\bookindexentry{词表转移}{Vocabulary Shifting}{\bookindexpage{section.6.3}{252}}
\bookindexentry{}{VoyageAI}{\bookindexpage{section.7.4}{283}}
\bookindexgroup{W}
\bookindexentry{Wasserstein 距离}{Wasserstein Distance}{\bookindexpage{section.6.2}{237}}
\bookindexentry{}{Whisper}{\bookindexpage{section.9.3}{379}}
\bookindexentry{白箱探查}{White-box Probing}{\bookindexpage{section.6.4}{258}}
\bookindexentry{}{WinoBias}{\bookindexpage{section.6.2}{237}, \bookindexpage{section.6.2}{244}}
\bookindexentry{}{Winogender}{\bookindexpage{section.6.2}{237}}
\bookindexentry{Winograd 模式挑战}{Winograd Schema Challenge}{\bookindexpage{section.6.2}{237}}
\bookindexentry{词嵌入关联检验}{Word Embedding Association Test}{\bookindexpage{section.6.2}{234}}
\bookindexentry{词嵌入}{Word Embeddings}{\bookindexpage{section.1.3}{9}}
\bookindexentry{词语过滤}{Word Filtering}{\bookindexpage{section.6.3}{252}}
\bookindexentry{}{word2vec}{\bookindexpage{section.1.3}{9}}
\bookindexentry{工作流编排}{Workflow Orchestration}{\bookindexpage{section.8.8}{354}}
\bookindexgroup{Y}
\bookindexentry{}{You.com}{\bookindexpage{section.10.4}{434}}
\bookindexgroup{Z}
\bookindexentry{零样本}{Zero-Shot}{\bookindexpage{section.2.5}{48}, \bookindexpage{section.3.1}{89}, \bookindexpage{section.4.1}{135}}
\bookindexentry{零样本提示}{Zero-Shot Prompts}{\bookindexpage{section.4.5}{163}}
\bookindexentry{}{ZeroGPU}{\bookindexpage{section.10.3}{430}}
\bookindexentry{}{ZeroQuant}{\bookindexpage{section.4.4}{159}}
\bookindexentry{}{Zeroscope}{\bookindexpage{section.9.6}{411}}

\end{multicols}
\endgroup
